
\documentclass{article}
\usepackage{colortbl}
\usepackage{makecell}
\usepackage{multirow}
\usepackage{supertabular}

\begin{document}

\newcounter{utterance}

\centering \large Interaction Transcript for game `matchit\_ascii', experiment `same\_grid', episode 3 with Qwen3.6{-}35B{-}A3B{-}FP8{-}without{-}reasoning.
\vspace{24pt}

{ \footnotesize  \setcounter{utterance}{1}
\setlength{\tabcolsep}{0pt}
\begin{supertabular}{c@{$\;$}|p{.15\linewidth}@{}p{.15\linewidth}p{.15\linewidth}p{.15\linewidth}p{.15\linewidth}p{.15\linewidth}}
    \# & $\;$A & \multicolumn{4}{c}{Game Master} & $\;\:$B\\
    \hline

    \theutterance \stepcounter{utterance}  
    & & \multicolumn{4}{p{0.6\linewidth}}{
        \cellcolor[rgb]{0.9,0.9,0.9}{
            \makecell[{{p{\linewidth}}}]{
                \texttt{\tiny{[P1$\langle$GM]}}
                \texttt{Dalyvaujate bendrame spėjimo žaidime.} \\
\texttt{Tikslas {-} išsiaiškinti, ar šis tinklelis ir kitas paveikslėlis, kurį matau tik aš, yra vienodi.} \\
\texttt{Pirmiausia apibūdinkite savo tinklelį. Tada aš pateiksiu savo aprašymą ir galėsime vienas kitam užduoti klausimų apie tinklelius, kad išsiaiškintume, ar jie vienodi. Dabar pradėkite savo trumpą tinklelio aprašymą žodžiais "APRAŠYMAS:", po kurių rašykite aprašymą. Nieko daugiau nepridėkite. Tai yra jūsų tinklelis:} \\
\texttt{X X X X X} \\
\texttt{▢ ▢ ▢ ▢ X} \\
\texttt{▢ ▢ X X X} \\
\texttt{▢ ▢ ▢ ▢ X} \\
\texttt{▢ ▢ ▢ ▢ X} \\
            }
        }
    }
    & \\ \\

    \theutterance \stepcounter{utterance}  
    & \multicolumn{4}{p{0.6\linewidth}}{
        \cellcolor[rgb]{0.9,0.9,0.9}{
            \makecell[{{p{\linewidth}}}]{
                \texttt{\tiny{[P1$\rangle$GM]}}
                \texttt{APRAŠYMAS: Tinklelyje yra 5 eilutės ir 5 stulpeliai. Pirmoje eilutėje visos 5 pozicijos užimtos X. Antroje eilutėje pirmosios 4 pozicijos yra tuščios (▢), o paskutinė – X. Trečioje eilutėje pirmosios dvi pozicijos yra tuščios, o paskutinės trys – X. Ketvirtoje ir penktoje eilutėse pirmosios keturios pozicijos yra tuščios, o paskutinė – X.} \\
            }
        }
    }
    & & \\ \\

    \theutterance \stepcounter{utterance}  
    & & & \multicolumn{2}{p{0.3\linewidth}}{
        \cellcolor[rgb]{0.9,0.9,0.9}{
            \makecell[{{p{\linewidth}}}]{
                \texttt{\tiny{[GM$|$GM]}}
                \texttt{continue} \\
            }
        }
    }
    & & \\ \\

    \theutterance \stepcounter{utterance}  
    & & \multicolumn{4}{p{0.6\linewidth}}{
        \cellcolor[rgb]{0.9,0.9,0.9}{
            \makecell[{{p{\linewidth}}}]{
                \texttt{\tiny{[GM$\rangle$P2]}}
                \texttt{Dalyvaujate bendrame spėjimo žaidime.} \\
\texttt{Tikslas {-} išsiaiškinti, ar šis tinklelis ir kitas paveikslėlis, kurį matau tik aš, yra vienodi.} \\
\texttt{Pirmiausia apibūdinkite savo tinklelį. Tada aš pateiksiu savo aprašymą ir galėsime vienas kitam užduoti klausimų apie tinklelius, kad išsiaiškintume, ar jie vienodi. Dabar pradėkite savo trumpą tinklelio aprašymą žodžiais "APRAŠYMAS:", po kurių rašykite aprašymą. Nieko daugiau nepridėkite. Tai yra jūsų tinklelis:} \\
\texttt{X X X X X} \\
\texttt{▢ ▢ ▢ ▢ X} \\
\texttt{▢ ▢ X X X} \\
\texttt{▢ ▢ ▢ ▢ X} \\
\texttt{▢ ▢ ▢ ▢ X} \\
            }
        }
    }
    & \\ \\

    \theutterance \stepcounter{utterance}  
    & & & \multicolumn{4}{p{0.6\linewidth}}{
        \cellcolor[rgb]{0.9,0.9,0.9}{
            \makecell[{{p{\linewidth}}}]{
                \texttt{\tiny{[GM$\langle$P2]}}
                \texttt{APRAŠYMAS: Tai 5x5 tinklelis, kuriame yra du skirtingi simboliai: "X" ir "▢". Tinklelyje yra 12 "X" ir 13 "▢". "X" simboliai sudaro specifinę figūrą: viršutinėje eilutėje visi 5 langeliai yra "X"; antrame stulpelyje nuo apačios iki trečios eilutės yra "X"; trečiame stulpelyje nuo trečios eilutės iki apačios yra "X"; ketvirtame stulpelyje nuo trečios eilutės iki apačios yra "X"; penktame stulpelyje antra, trečia, ketvirta ir penkta eilutės yra "X". Paprasčiau tariant, "X" užima visą pirmąją eilutę, o dešinėje pusėje sudaro vertikalią liniją penktame stulpelyje (išskyrus antrą eilutę, kuri jau užimta pirmos eilutės "X") ir horizontalią jungtį trečioje eilutėje. Tiksliau: (0,0){-}(0,4) yra X; (1,4) yra X; (2,2){-}(2,4) yra X; (3,4) yra X; (4,4) yra X.} \\
            }
        }
    }
     \\ \\

    \theutterance \stepcounter{utterance}  
    & & & \multicolumn{2}{p{0.3\linewidth}}{
        \cellcolor[rgb]{0.9,0.9,0.9}{
            \makecell[{{p{\linewidth}}}]{
                \texttt{\tiny{[GM$|$GM]}}
                \texttt{continue} \\
            }
        }
    }
    & & \\ \\

    \theutterance \stepcounter{utterance}  
    & & \multicolumn{4}{p{0.6\linewidth}}{
        \cellcolor[rgb]{0.9,0.9,0.9}{
            \makecell[{{p{\linewidth}}}]{
                \texttt{\tiny{[GM$\rangle$P2]}}
                \texttt{Tai mano APRAŠYMASAPRAŠYMAS: Tinklelyje yra 5 eilutės ir 5 stulpeliai. Pirmoje eilutėje visos 5 pozicijos užimtos X. Antroje eilutėje pirmosios 4 pozicijos yra tuščios (▢), o paskutinė – X. Trečioje eilutėje pirmosios dvi pozicijos yra tuščios, o paskutinės trys – X. Ketvirtoje ir penktoje eilutėse pirmosios keturios pozicijos yra tuščios, o paskutinė – X.} \\
\texttt{Dabar užduokite klausimą, kad sužinotumėte naujus mano tinklelio aspektus, kurie gali skirtis nuo jūsų tinklelio. Pradėkite "KLAUSIMAS:" ir nieko daugiau nepridėkite.} \\
            }
        }
    }
    & \\ \\

    \theutterance \stepcounter{utterance}  
    & & & \multicolumn{4}{p{0.6\linewidth}}{
        \cellcolor[rgb]{0.9,0.9,0.9}{
            \makecell[{{p{\linewidth}}}]{
                \texttt{\tiny{[GM$\langle$P2]}}
                \texttt{KLAUSIMAS: Ar jūsų tinklelyje trečioje eilutėje trečiasis stulpelis (vidurinis) yra užpildytas simboliu "X"?} \\
            }
        }
    }
     \\ \\

    \theutterance \stepcounter{utterance}  
    & & & \multicolumn{2}{p{0.3\linewidth}}{
        \cellcolor[rgb]{0.9,0.9,0.9}{
            \makecell[{{p{\linewidth}}}]{
                \texttt{\tiny{[GM$|$GM]}}
                \texttt{continue} \\
            }
        }
    }
    & & \\ \\

    \theutterance \stepcounter{utterance}  
    & & \multicolumn{4}{p{0.6\linewidth}}{
        \cellcolor[rgb]{0.9,0.9,0.9}{
            \makecell[{{p{\linewidth}}}]{
                \texttt{\tiny{[P1$\langle$GM]}}
                \texttt{Tai mano APRAŠYMASAPRAŠYMAS: Tai 5x5 tinklelis, kuriame yra du skirtingi simboliai: "X" ir "▢". Tinklelyje yra 12 "X" ir 13 "▢". "X" simboliai sudaro specifinę figūrą: viršutinėje eilutėje visi 5 langeliai yra "X"; antrame stulpelyje nuo apačios iki trečios eilutės yra "X"; trečiame stulpelyje nuo trečios eilutės iki apačios yra "X"; ketvirtame stulpelyje nuo trečios eilutės iki apačios yra "X"; penktame stulpelyje antra, trečia, ketvirta ir penkta eilutės yra "X". Paprasčiau tariant, "X" užima visą pirmąją eilutę, o dešinėje pusėje sudaro vertikalią liniją penktame stulpelyje (išskyrus antrą eilutę, kuri jau užimta pirmos eilutės "X") ir horizontalią jungtį trečioje eilutėje. Tiksliau: (0,0){-}(0,4) yra X; (1,4) yra X; (2,2){-}(2,4) yra X; (3,4) yra X; (4,4) yra X.} \\
\texttt{KLAUSIMAS: Ar jūsų tinklelyje trečioje eilutėje trečiasis stulpelis (vidurinis) yra užpildytas simboliu "X"?Pradėkite atsakymą „ATSAKYMAS:" ir nieko daugiau nepridėkite.} \\
            }
        }
    }
    & \\ \\

    \theutterance \stepcounter{utterance}  
    & \multicolumn{4}{p{0.6\linewidth}}{
        \cellcolor[rgb]{0.9,0.9,0.9}{
            \makecell[{{p{\linewidth}}}]{
                \texttt{\tiny{[P1$\rangle$GM]}}
                \texttt{ATSAKYMAS: Taip, mano tinklelyje trečioje eilutėje trečiasis stulpelis yra užpildytas simboliu "X".} \\
            }
        }
    }
    & & \\ \\

    \theutterance \stepcounter{utterance}  
    & & & \multicolumn{2}{p{0.3\linewidth}}{
        \cellcolor[rgb]{0.9,0.9,0.9}{
            \makecell[{{p{\linewidth}}}]{
                \texttt{\tiny{[GM$|$GM]}}
                \texttt{continue} \\
            }
        }
    }
    & & \\ \\

    \theutterance \stepcounter{utterance}  
    & & \multicolumn{4}{p{0.6\linewidth}}{
        \cellcolor[rgb]{0.9,0.9,0.9}{
            \makecell[{{p{\linewidth}}}]{
                \texttt{\tiny{[P1$\langle$GM]}}
                \texttt{Dabar užduokite klausimą, kad sužinotumėte naujus mano tinklelio aspektus, kurie gali skirtis nuo jūsų tinklelio. Pradėkite "KLAUSIMAS:" ir nieko daugiau nepridėkite.} \\
            }
        }
    }
    & \\ \\

    \theutterance \stepcounter{utterance}  
    & \multicolumn{4}{p{0.6\linewidth}}{
        \cellcolor[rgb]{0.9,0.9,0.9}{
            \makecell[{{p{\linewidth}}}]{
                \texttt{\tiny{[P1$\rangle$GM]}}
                \texttt{KLAUSIMAS: Ar jūsų tinklelyje antroje eilutėje pirmasis stulpelis yra užpildytas simboliu "X"?} \\
            }
        }
    }
    & & \\ \\

    \theutterance \stepcounter{utterance}  
    & & & \multicolumn{2}{p{0.3\linewidth}}{
        \cellcolor[rgb]{0.9,0.9,0.9}{
            \makecell[{{p{\linewidth}}}]{
                \texttt{\tiny{[GM$|$GM]}}
                \texttt{continue} \\
            }
        }
    }
    & & \\ \\

    \theutterance \stepcounter{utterance}  
    & & \multicolumn{4}{p{0.6\linewidth}}{
        \cellcolor[rgb]{0.9,0.9,0.9}{
            \makecell[{{p{\linewidth}}}]{
                \texttt{\tiny{[GM$\rangle$P2]}}
                \texttt{ATSAKYMAS: Taip, mano tinklelyje trečioje eilutėje trečiasis stulpelis yra užpildytas simboliu "X".} \\
\texttt{KLAUSIMAS: Ar jūsų tinklelyje antroje eilutėje pirmasis stulpelis yra užpildytas simboliu "X"?Pradėkite atsakymą „ATSAKYMAS:" ir nieko daugiau nepridėkite.} \\
            }
        }
    }
    & \\ \\

    \theutterance \stepcounter{utterance}  
    & & & \multicolumn{4}{p{0.6\linewidth}}{
        \cellcolor[rgb]{0.9,0.9,0.9}{
            \makecell[{{p{\linewidth}}}]{
                \texttt{\tiny{[GM$\langle$P2]}}
                \texttt{ATSAKYMAS: Ne, mano tinklelyje antroje eilutėje pirmasis stulpelis yra užpildytas simboliu "▢".} \\
            }
        }
    }
     \\ \\

    \theutterance \stepcounter{utterance}  
    & & & \multicolumn{2}{p{0.3\linewidth}}{
        \cellcolor[rgb]{0.9,0.9,0.9}{
            \makecell[{{p{\linewidth}}}]{
                \texttt{\tiny{[GM$|$GM]}}
                \texttt{continue} \\
            }
        }
    }
    & & \\ \\

    \theutterance \stepcounter{utterance}  
    & & \multicolumn{4}{p{0.6\linewidth}}{
        \cellcolor[rgb]{0.9,0.9,0.9}{
            \makecell[{{p{\linewidth}}}]{
                \texttt{\tiny{[GM$\rangle$P2]}}
                \texttt{Dabar užduokite klausimą, kad sužinotumėte naujus mano tinklelio aspektus, kurie gali skirtis nuo jūsų tinklelio. Pradėkite "KLAUSIMAS:" ir nieko daugiau nepridėkite.} \\
            }
        }
    }
    & \\ \\

    \theutterance \stepcounter{utterance}  
    & & & \multicolumn{4}{p{0.6\linewidth}}{
        \cellcolor[rgb]{0.9,0.9,0.9}{
            \makecell[{{p{\linewidth}}}]{
                \texttt{\tiny{[GM$\langle$P2]}}
                \texttt{KLAUSIMAS: Ar jūsų tinklelyje ketvirtoje eilutėje trečiasis stulpelis yra užpildytas simboliu "X"?} \\
            }
        }
    }
     \\ \\

    \theutterance \stepcounter{utterance}  
    & & & \multicolumn{2}{p{0.3\linewidth}}{
        \cellcolor[rgb]{0.9,0.9,0.9}{
            \makecell[{{p{\linewidth}}}]{
                \texttt{\tiny{[GM$|$GM]}}
                \texttt{continue} \\
            }
        }
    }
    & & \\ \\

    \theutterance \stepcounter{utterance}  
    & & \multicolumn{4}{p{0.6\linewidth}}{
        \cellcolor[rgb]{0.9,0.9,0.9}{
            \makecell[{{p{\linewidth}}}]{
                \texttt{\tiny{[P1$\langle$GM]}}
                \texttt{ATSAKYMAS: Ne, mano tinklelyje antroje eilutėje pirmasis stulpelis yra užpildytas simboliu "▢".} \\
\texttt{KLAUSIMAS: Ar jūsų tinklelyje ketvirtoje eilutėje trečiasis stulpelis yra užpildytas simboliu "X"?Pradėkite atsakymą „ATSAKYMAS:" ir nieko daugiau nepridėkite.} \\
            }
        }
    }
    & \\ \\

    \theutterance \stepcounter{utterance}  
    & \multicolumn{4}{p{0.6\linewidth}}{
        \cellcolor[rgb]{0.9,0.9,0.9}{
            \makecell[{{p{\linewidth}}}]{
                \texttt{\tiny{[P1$\rangle$GM]}}
                \texttt{ATSAKYMAS: Ne, mano tinklelyje ketvirtoje eilutėje trečiasis stulpelis yra užpildytas simboliu "▢".} \\
            }
        }
    }
    & & \\ \\

    \theutterance \stepcounter{utterance}  
    & & & \multicolumn{2}{p{0.3\linewidth}}{
        \cellcolor[rgb]{0.9,0.9,0.9}{
            \makecell[{{p{\linewidth}}}]{
                \texttt{\tiny{[GM$|$GM]}}
                \texttt{continue} \\
            }
        }
    }
    & & \\ \\

    \theutterance \stepcounter{utterance}  
    & & \multicolumn{4}{p{0.6\linewidth}}{
        \cellcolor[rgb]{0.9,0.9,0.9}{
            \makecell[{{p{\linewidth}}}]{
                \texttt{\tiny{[P1$\langle$GM]}}
                \texttt{Dabar užduokite klausimą, kad sužinotumėte naujus mano tinklelio aspektus, kurie gali skirtis nuo jūsų tinklelio. Pradėkite "KLAUSIMAS:" ir nieko daugiau nepridėkite.} \\
            }
        }
    }
    & \\ \\

    \theutterance \stepcounter{utterance}  
    & \multicolumn{4}{p{0.6\linewidth}}{
        \cellcolor[rgb]{0.9,0.9,0.9}{
            \makecell[{{p{\linewidth}}}]{
                \texttt{\tiny{[P1$\rangle$GM]}}
                \texttt{KLAUSIMAS: Ar jūsų tinklelyje trečioje eilutėje pirmasis stulpelis yra užpildytas simboliu "X"?} \\
            }
        }
    }
    & & \\ \\

    \theutterance \stepcounter{utterance}  
    & & & \multicolumn{2}{p{0.3\linewidth}}{
        \cellcolor[rgb]{0.9,0.9,0.9}{
            \makecell[{{p{\linewidth}}}]{
                \texttt{\tiny{[GM$|$GM]}}
                \texttt{continue} \\
            }
        }
    }
    & & \\ \\

    \theutterance \stepcounter{utterance}  
    & & \multicolumn{4}{p{0.6\linewidth}}{
        \cellcolor[rgb]{0.9,0.9,0.9}{
            \makecell[{{p{\linewidth}}}]{
                \texttt{\tiny{[GM$\rangle$P2]}}
                \texttt{ATSAKYMAS: Ne, mano tinklelyje ketvirtoje eilutėje trečiasis stulpelis yra užpildytas simboliu "▢".} \\
\texttt{KLAUSIMAS: Ar jūsų tinklelyje trečioje eilutėje pirmasis stulpelis yra užpildytas simboliu "X"?Pradėkite atsakymą „ATSAKYMAS:" ir nieko daugiau nepridėkite.} \\
            }
        }
    }
    & \\ \\

    \theutterance \stepcounter{utterance}  
    & & & \multicolumn{4}{p{0.6\linewidth}}{
        \cellcolor[rgb]{0.9,0.9,0.9}{
            \makecell[{{p{\linewidth}}}]{
                \texttt{\tiny{[GM$\langle$P2]}}
                \texttt{ATSAKYMAS: Ne, mano tinklelyje trečioje eilutėje pirmasis stulpelis yra užpildytas simboliu "▢".} \\
            }
        }
    }
     \\ \\

    \theutterance \stepcounter{utterance}  
    & & & \multicolumn{2}{p{0.3\linewidth}}{
        \cellcolor[rgb]{0.9,0.9,0.9}{
            \makecell[{{p{\linewidth}}}]{
                \texttt{\tiny{[GM$|$GM]}}
                \texttt{continue} \\
            }
        }
    }
    & & \\ \\

    \theutterance \stepcounter{utterance}  
    & & \multicolumn{4}{p{0.6\linewidth}}{
        \cellcolor[rgb]{0.9,0.9,0.9}{
            \makecell[{{p{\linewidth}}}]{
                \texttt{\tiny{[GM$\rangle$P2]}}
                \texttt{Dabar užduokite klausimą, kad sužinotumėte naujus mano tinklelio aspektus, kurie gali skirtis nuo jūsų tinklelio. Pradėkite "KLAUSIMAS:" ir nieko daugiau nepridėkite.} \\
            }
        }
    }
    & \\ \\

    \theutterance \stepcounter{utterance}  
    & & & \multicolumn{4}{p{0.6\linewidth}}{
        \cellcolor[rgb]{0.9,0.9,0.9}{
            \makecell[{{p{\linewidth}}}]{
                \texttt{\tiny{[GM$\langle$P2]}}
                \texttt{KLAUSIMAS: Ar jūsų tinklelyje penktoje eilutėje antrasis stulpelis yra užpildytas simboliu "X"?} \\
            }
        }
    }
     \\ \\

    \theutterance \stepcounter{utterance}  
    & & & \multicolumn{2}{p{0.3\linewidth}}{
        \cellcolor[rgb]{0.9,0.9,0.9}{
            \makecell[{{p{\linewidth}}}]{
                \texttt{\tiny{[GM$|$GM]}}
                \texttt{continue} \\
            }
        }
    }
    & & \\ \\

    \theutterance \stepcounter{utterance}  
    & & \multicolumn{4}{p{0.6\linewidth}}{
        \cellcolor[rgb]{0.9,0.9,0.9}{
            \makecell[{{p{\linewidth}}}]{
                \texttt{\tiny{[P1$\langle$GM]}}
                \texttt{ATSAKYMAS: Ne, mano tinklelyje trečioje eilutėje pirmasis stulpelis yra užpildytas simboliu "▢".} \\
\texttt{KLAUSIMAS: Ar jūsų tinklelyje penktoje eilutėje antrasis stulpelis yra užpildytas simboliu "X"?Pradėkite atsakymą „ATSAKYMAS:" ir nieko daugiau nepridėkite.} \\
            }
        }
    }
    & \\ \\

    \theutterance \stepcounter{utterance}  
    & \multicolumn{4}{p{0.6\linewidth}}{
        \cellcolor[rgb]{0.9,0.9,0.9}{
            \makecell[{{p{\linewidth}}}]{
                \texttt{\tiny{[P1$\rangle$GM]}}
                \texttt{ATSAKYMAS: Ne, mano tinklelyje penktoje eilutėje antrasis stulpelis yra užpildytas simboliu "▢".} \\
            }
        }
    }
    & & \\ \\

    \theutterance \stepcounter{utterance}  
    & & & \multicolumn{2}{p{0.3\linewidth}}{
        \cellcolor[rgb]{0.9,0.9,0.9}{
            \makecell[{{p{\linewidth}}}]{
                \texttt{\tiny{[GM$|$GM]}}
                \texttt{continue} \\
            }
        }
    }
    & & \\ \\

    \theutterance \stepcounter{utterance}  
    & & \multicolumn{4}{p{0.6\linewidth}}{
        \cellcolor[rgb]{0.9,0.9,0.9}{
            \makecell[{{p{\linewidth}}}]{
                \texttt{\tiny{[P1$\langle$GM]}}
                \texttt{Dabar užduokite klausimą, kad sužinotumėte naujus mano tinklelio aspektus, kurie gali skirtis nuo jūsų tinklelio. Pradėkite "KLAUSIMAS:" ir nieko daugiau nepridėkite.} \\
            }
        }
    }
    & \\ \\

    \theutterance \stepcounter{utterance}  
    & \multicolumn{4}{p{0.6\linewidth}}{
        \cellcolor[rgb]{0.9,0.9,0.9}{
            \makecell[{{p{\linewidth}}}]{
                \texttt{\tiny{[P1$\rangle$GM]}}
                \texttt{KLAUSIMAS: Ar jūsų tinklelyje antroje eilutėje trečiasis stulpelis yra užpildytas simboliu "X"?} \\
            }
        }
    }
    & & \\ \\

    \theutterance \stepcounter{utterance}  
    & & & \multicolumn{2}{p{0.3\linewidth}}{
        \cellcolor[rgb]{0.9,0.9,0.9}{
            \makecell[{{p{\linewidth}}}]{
                \texttt{\tiny{[GM$|$GM]}}
                \texttt{continue} \\
            }
        }
    }
    & & \\ \\

    \theutterance \stepcounter{utterance}  
    & & \multicolumn{4}{p{0.6\linewidth}}{
        \cellcolor[rgb]{0.9,0.9,0.9}{
            \makecell[{{p{\linewidth}}}]{
                \texttt{\tiny{[GM$\rangle$P2]}}
                \texttt{ATSAKYMAS: Ne, mano tinklelyje penktoje eilutėje antrasis stulpelis yra užpildytas simboliu "▢".} \\
\texttt{KLAUSIMAS: Ar jūsų tinklelyje antroje eilutėje trečiasis stulpelis yra užpildytas simboliu "X"?Pradėkite atsakymą „ATSAKYMAS:" ir nieko daugiau nepridėkite.} \\
            }
        }
    }
    & \\ \\

    \theutterance \stepcounter{utterance}  
    & & & \multicolumn{4}{p{0.6\linewidth}}{
        \cellcolor[rgb]{0.9,0.9,0.9}{
            \makecell[{{p{\linewidth}}}]{
                \texttt{\tiny{[GM$\langle$P2]}}
                \texttt{ATSAKYMAS: Ne, mano tinklelyje antroje eilutėje trečiasis stulpelis yra užpildytas simboliu "▢".} \\
            }
        }
    }
     \\ \\

    \theutterance \stepcounter{utterance}  
    & & & \multicolumn{2}{p{0.3\linewidth}}{
        \cellcolor[rgb]{0.9,0.9,0.9}{
            \makecell[{{p{\linewidth}}}]{
                \texttt{\tiny{[GM$|$GM]}}
                \texttt{continue} \\
            }
        }
    }
    & & \\ \\

    \theutterance \stepcounter{utterance}  
    & & \multicolumn{4}{p{0.6\linewidth}}{
        \cellcolor[rgb]{0.9,0.9,0.9}{
            \makecell[{{p{\linewidth}}}]{
                \texttt{\tiny{[GM$\rangle$P2]}}
                \texttt{Dabar priimkite sprendimą. Kaip manote: ar jūsų tinklelis ir kitas aprašytas tinklelis yra tokie patys? Parašykite "SPRENDIMAS: tas pats tinklelis", jei manote, kad tai tas pats tinklelis, arba "SPRENDIMAS: skirtingi tinkleliai", jei manote, kad tai skirtingi tinkleliai. Nieko daugiau nepridėkite.} \\
            }
        }
    }
    & \\ \\

    \theutterance \stepcounter{utterance}  
    & & & \multicolumn{4}{p{0.6\linewidth}}{
        \cellcolor[rgb]{0.9,0.9,0.9}{
            \makecell[{{p{\linewidth}}}]{
                \texttt{\tiny{[GM$\langle$P2]}}
                \texttt{SPRENDIMAS: tas pats tinklelis} \\
            }
        }
    }
     \\ \\

    \theutterance \stepcounter{utterance}  
    & & & \multicolumn{2}{p{0.3\linewidth}}{
        \cellcolor[rgb]{0.9,0.9,0.9}{
            \makecell[{{p{\linewidth}}}]{
                \texttt{\tiny{[GM$|$GM]}}
                \texttt{continue} \\
            }
        }
    }
    & & \\ \\

    \theutterance \stepcounter{utterance}  
    & & & \multicolumn{2}{p{0.3\linewidth}}{
        \cellcolor[rgb]{0.9,0.9,0.9}{
            \makecell[{{p{\linewidth}}}]{
                \texttt{\tiny{[GM$|$GM]}}
                \texttt{success} \\
            }
        }
    }
    & & \\ \\

    \theutterance \stepcounter{utterance}  
    & & \multicolumn{4}{p{0.6\linewidth}}{
        \cellcolor[rgb]{0.9,0.9,0.9}{
            \makecell[{{p{\linewidth}}}]{
                \texttt{\tiny{[P1$\langle$GM]}}
                \texttt{ATSAKYMAS: Ne, mano tinklelyje antroje eilutėje trečiasis stulpelis yra užpildytas simboliu "▢".} \\
\texttt{Dabar priimkite sprendimą. Kaip manote: ar jūsų tinklelis ir kitas aprašytas tinklelis yra tokie patys? Parašykite "SPRENDIMAS: tas pats tinklelis", jei manote, kad tai tas pats tinklelis, arba "SPRENDIMAS: skirtingi tinkleliai", jei manote, kad tai skirtingi tinkleliai. Nieko daugiau nepridėkite.} \\
            }
        }
    }
    & \\ \\

    \theutterance \stepcounter{utterance}  
    & \multicolumn{4}{p{0.6\linewidth}}{
        \cellcolor[rgb]{0.9,0.9,0.9}{
            \makecell[{{p{\linewidth}}}]{
                \texttt{\tiny{[P1$\rangle$GM]}}
                \texttt{SPRENDIMAS: tas pats tinklelis} \\
            }
        }
    }
    & & \\ \\

    \theutterance \stepcounter{utterance}  
    & & & \multicolumn{2}{p{0.3\linewidth}}{
        \cellcolor[rgb]{0.9,0.9,0.9}{
            \makecell[{{p{\linewidth}}}]{
                \texttt{\tiny{[GM$|$GM]}}
                \texttt{continue} \\
            }
        }
    }
    & & \\ \\

    \theutterance \stepcounter{utterance}  
    & & & \multicolumn{2}{p{0.3\linewidth}}{
        \cellcolor[rgb]{0.9,0.9,0.9}{
            \makecell[{{p{\linewidth}}}]{
                \texttt{\tiny{[GM$|$GM]}}
                \texttt{success} \\
            }
        }
    }
    & & \\ \\

\end{supertabular}
}

\end{document}
